\section*{Chapter \ref{ch:g10:the-mole} --- The Mole and Molar Mass}

\begin{solution}{exo:g10:the-mole:1}
$M(\ce{O2}) = 2 \times 16.0 = \qty{32.0}{g/mol}$;
$M(\ce{N2}) = 2 \times 14.0 = \qty{28.0}{g/mol}$;
$M(\ce{CH4}) = 12.0 + 4 \times 1.0 = \qty{16.0}{g/mol}$;
$M(\ce{NH3}) = 14.0 + 3 \times 1.0 = \qty{17.0}{g/mol}$;
$M(\ce{CO2}) = 12.0 + 2 \times 16.0 = \qty{44.0}{g/mol}$.
\end{solution}

\begin{solution}{exo:g10:the-mole:2}
Water: $n = 36.0 / 18.0 = \qty{2.00}{mol}$. Methane:
$n = 4.0 / 16.0 = \qty{0.25}{mol}$.
\end{solution}

\begin{solution}{exo:g10:the-mole:3}
$m(\ce{NaCl}) = 0.25 \times 58.5 = \qty{14.6}{g}$, about \qty{15}{g};
$m(\ce{CO2}) = 2.0 \times 44.0 = \qty{88}{g}$.
\end{solution}

\begin{solution}{exo:g10:the-mole:4}
$N = 0.50 \times \num{6.02e23} = \num{3.0e23}$ \omterm{def:g7:atoms-and-molecules:molecule}{molecules} of dioxygen.
$n = \num{3.0e24} / \qty{6.02e23}{mol^{-1}} = \qty{5.0}{mol}$ of iron.
\end{solution}

\begin{solution}{exo:g10:the-mole:5}
$V = n \times V_m = 0.50 \times 24.1 = \qty{12}{L}$.
$n = V / V_m = 12 / 24.1 = \qty{0.50}{mol}$.
\end{solution}

\begin{solution}{exo:g10:the-mole:6}
The drop has a mass of \qty{0.050}{g}, so
$n = 0.050 / 18.0 = \qty{2.8e-3}{mol}$ and
$N = \num{2.8e-3} \times \num{6.02e23} = \num{1.7e21}$ \omterm{def:g7:atoms-and-molecules:molecule}{molecules}.
\end{solution}

\begin{solution}{exo:g10:the-mole:7}
An aluminium \omterm{def:g7:atoms-and-molecules:atom}{atom} (\qty{27.0}{g/mol}) is about half as heavy as an iron
\omterm{def:g7:atoms-and-molecules:atom}{atom} (\qty{55.8}{g/mol}), so the same mass of aluminium holds about
twice as many \omterm{def:g7:atoms-and-molecules:atom}{atoms}. Check: aluminium $10/27.0 = \qty{0.370}{mol}$, that
is \num{2.23e23} \omterm{def:g7:atoms-and-molecules:atom}{atoms}; iron $10/55.8 = \qty{0.179}{mol}$, that is
\num{1.08e23} \omterm{def:g7:atoms-and-molecules:atom}{atoms}. Aluminium wins, by a factor of about 2.1.
\end{solution}

\begin{solution}{exo:g10:the-mole:8}
One litre is $1/24.1 = \qty{0.0415}{mol}$ of gas. Carbon dioxide:
$0.0415 \times 44.0 = \qty{1.83}{g}$; dinitrogen:
$0.0415 \times 28.0 = \qty{1.16}{g}$. Carbon dioxide is about one and a
half times denser than \omterm{def:g4:air-a-mixture-of-gases:air}{air}, whose \omterm{def:g10:the-mole:molar-mass}{molar mass} is close to that of
dinitrogen. The fermenting juice gives off carbon dioxide, which sinks
and builds up near the floor of a closed cellar, where it can suffocate
anyone who goes in.
\end{solution}

\begin{solution}{exo:g10:the-mole:9}
$M(\ce{C6H12O6}) = 6 \times 12.0 + 12 \times 1.0 + 6 \times 16.0 =
\qty{180.0}{g/mol}$. Arrow $\div N_A$: $n = \num{1.0e22} / \num{6.02e23} =
\qty{1.66e-2}{mol}$; arrow $\times M$: $m = \num{1.66e-2} \times 180.0 =
\qty{3.0}{g}$.
\end{solution}

\begin{solution}{exo:g10:the-mole:10}
$M = \qty{342.0}{g/mol}$, so $n = 5.0 / 342.0 = \qty{1.46e-2}{mol}$;
$N = \num{1.46e-2} \times \num{6.02e23} = \num{8.8e21}$ \omterm{def:g7:atoms-and-molecules:molecule}{molecules}. Each
has 12 carbon \omterm{def:g7:atoms-and-molecules:atom}{atoms}: $12 \times \num{8.8e21} = \num{1.1e23}$ carbon \omterm{def:g7:atoms-and-molecules:atom}{atoms}.
\end{solution}

\begin{solution}{exo:g10:the-mole:11}
Water: $1000 / 18.0 = \qty{55.6}{mol}$. Ethanol:
$M = 2 \times 12.0 + 6 \times 1.0 + 16.0 = \qty{46.0}{g/mol}$, so
$1000 / 46.0 = \qty{21.7}{mol}$. Water has the larger amount, by a
factor $55.6 / 21.7 \approx 2.6$, which is simply $46.0 / 18.0$.
\end{solution}

\begin{solution}{exo:g10:the-mole:12}
Gold: $\qty{3.75}{g}$, $n = 3.75 / 197.0 = \qty{1.90e-2}{mol}$, that is
\num{1.15e22} \omterm{def:g7:atoms-and-molecules:atom}{atoms}. Copper: \qty{1.25}{g},
$n = 1.25 / 63.5 = \qty{1.97e-2}{mol}$, that is \num{1.19e22} \omterm{def:g7:atoms-and-molecules:atom}{atoms}.
Surprisingly, the ring holds slightly more copper \omterm{def:g7:atoms-and-molecules:atom}{atoms} than gold \omterm{def:g7:atoms-and-molecules:atom}{atoms}:
a gold \omterm{def:g7:atoms-and-molecules:atom}{atom} is about three times heavier than a copper \omterm{def:g7:atoms-and-molecules:atom}{atom}.
\end{solution}

\begin{solution}{exo:g10:the-mole:13}
\begin{enumerate}
  \item $V = 8.0 \times 6.0 \times 3.0 = \qty{144}{m^3} =
        \qty{1.44e5}{L}$; $n = \num{1.44e5} / 24.1 = \qty{5.98e3}{mol}$;
        $N = \num{5.98e3} \times \num{6.02e23} = \num{3.6e27}$ \omterm{def:g7:atoms-and-molecules:molecule}{molecules}.
  \item $M = (78 \times 28.0 + 21 \times 32.0 + 1 \times 40.0) / 100 =
        2896 / 100 = \qty{29.0}{g/mol}$.
  \item $m = \num{5.98e3} \times 29.0 = \qty{1.73e5}{g}$, about
        \qty{173}{kg}: the \omterm{def:g4:air-a-mixture-of-gases:air}{air} of a classroom weighs as much as two or
        three adults.
\end{enumerate}
\end{solution}

\begin{solution}{exo:g10:the-mole:14}
\begin{enumerate}
  \item $n = 1000 / 28.1 = \qty{35.6}{mol}$, so
        $N = 35.6 \times \num{6.02e23} = \num{2.14e25}$ \omterm{def:g7:atoms-and-molecules:atom}{atoms}.
  \item $m = \qty{1.000}{kg} / \num{2.14e25} = \qty{4.67e-26}{kg}$.
  \item The mass of the sphere and the \omterm{def:g10:the-mole:molar-mass}{molar mass} give its amount,
        $n = m / M$. If the \omterm{def:g7:atoms-and-molecules:atom}{atoms} $N$ are counted independently (from
        the volume of the sphere and the spacing of the \omterm{def:g7:atoms-and-molecules:atom}{atoms} in the
        crystal), the ratio $N / n$ is the \omterm{def:g10:the-mole:avogadro}{Avogadro constant}.
\end{enumerate}
\end{solution}

\begin{solution}{exo:g10:the-mole:15}
$M = 0.758 \times 35.0 + 0.242 \times 37.0 = 26.53 + 8.95 =
\qty{35.5}{g/mol}$: the value of the table.
\end{solution}

\begin{solution}{pb:g10:the-mole:1}
\textbf{1.} $M(\ce{H2O}) = 2 \times 1.0 + 16.0 = \qty{18.0}{g/mol}$.

\textbf{2.} $m = 250 \times \qty{1.0}{g} = \qty{250}{g}$.

\textbf{3.} $n = 250 / 18.0 = \qty{13.9}{mol}$.

\textbf{4.} Each \omterm{def:g7:atoms-and-molecules:molecule}{molecule} has two hydrogen \omterm{def:g7:atoms-and-molecules:atom}{atoms} and one oxygen \omterm{def:g7:atoms-and-molecules:atom}{atom}:
\qty{27.8}{mol} of hydrogen \omterm{def:g7:atoms-and-molecules:atom}{atoms}, \qty{13.9}{mol} of oxygen \omterm{def:g7:atoms-and-molecules:atom}{atoms}.

\textbf{5.} $N = 13.9 \times \num{6.02e23} = \num{8.36e24}$ \omterm{def:g7:atoms-and-molecules:molecule}{molecules}.

\textbf{6.} $m = \qty{18.0}{g/mol} / \qty{6.02e23}{mol^{-1}} =
\qty{2.99e-23}{g}$.

\textbf{7.} $\num{8.36e24} \times \qty{2.99e-23}{g} = \qty{250}{g}$, the
mass of the water in the glass, as it should be.

\textbf{8.} $\num{8.36e24} / \num{3.16e7} = \num{2.6e17}$ years: a
number of years with seventeen zeros, far beyond any possible count.

\textbf{9.} $\qty{1.335e9}{km^3} \times \num{e9} = \qty{1.335e18}{m^3}$,
and with $\qty{1}{m^3} = \qty{1000}{L}$, \qty{1.335e21}{L}.

\textbf{10.} $\num{8.36e24} / \num{1.335e21} = \num{6.26e3}$ marked
\omterm{def:g7:atoms-and-molecules:molecule}{molecules} per litre.

\textbf{11.} $\num{6.26e3} \times 0.250 = \num{1.57e3}$ marked \omterm{def:g7:atoms-and-molecules:molecule}{molecules}
in the second glass.

\textbf{12.} Mass $\num{1.335e21} \times \qty{1.0}{kg} =
\qty{1.335e21}{kg} = \qty{1.335e24}{g}$; $n = \num{1.335e24} / 18.0 =
\qty{7.42e22}{mol}$.

\textbf{13.} $13.9 / \num{7.42e22} = \num{1.87e-22}$.

\textbf{14.} $\num{1.87e-22} \times \num{8.36e24} = \num{1.57e3}$: the
same answer as question 11, as it must be.

\textbf{15.} $\num{1.335e21} / 0.250 = \num{5.34e21}$ glasses.

\textbf{16.} $\num{8.36e24} / \num{5.34e21} \approx 1570$. A glass holds
about 1600 times more \omterm{def:g7:atoms-and-molecules:molecule}{molecules} than the oceans hold glasses of water;
spread evenly, the \omterm{def:g7:atoms-and-molecules:molecule}{molecules} of one glass leave about 1600 in every
glassful of ocean.

\textbf{17.} $1 / \num{6.26e3} = \qty{1.6e-4}{L} = \qty{0.16}{mL}$, about
three drops.

\textbf{18.} The ocean volume is an estimate known to three or four
figures at best, and it was rounded; sea water is not pure water (it
holds salts); the glass is not exactly \qty{250}{mL}; and the mixing
through all the oceans would never be perfectly even. Only the order of
magnitude can be trusted.

\textbf{19.} About 1600 \omterm{def:g7:atoms-and-molecules:molecule}{molecules} of the first glass come back in the
second.
\end{solution}
